\documentclass[11pt]{article}
\usepackage[T1]{fontenc}
\usepackage{lmodern,amsmath,amssymb,amsthm,booktabs,microtype}
\usepackage[margin=1in]{geometry}
\usepackage[colorlinks=true,urlcolor=blue,citecolor=blue]{hyperref}
\newtheorem{theorem}{Theorem}
\newtheorem{corollary}[theorem]{Corollary}
\newcommand{\A}[1]{a_{#1}}
\title{A classical star family and the affine-$D_4$ boundary}
\author{Explicit proof and independent computational certificate}
\date{6 October 2026}
\begin{document}
\maketitle
\begin{abstract}
A broad search rediscovered a simple fully commutative-bottom family with
Catalan leading coefficients. Its four-leaf case is affine $D_4$ and has
$P_{x,w}=1+3q+2q^2$. We give an elementary proof and an independent
affine-action computation. This supplies a precise obstruction to extending
the type-$D$ fully commutative-bottom bound to every Coxeter system. The
exact star family is explicitly covered by Mongelli's 2011 preprint,
Corollary 2, as well as older cube-interval formulas; it is not presented
as a new discovery or as a use requiring the
September 2026 invariance theorem.
\end{abstract}

\section{The star Coxeter systems}
Let $W_m$ have generators $c,a_1,\ldots,a_m$, with
$m(c,a_i)=3$ and $m(a_i,a_j)=2$ for $i\ne j$.
For $L\subseteq[m]$ put $a_L=\prod_{i\in L}a_i$ and
\[
 x=a_{[m]},\qquad w=xcx.
\]
The lower endpoint is a product of distinct commuting generators, hence
fully commutative. The systems for $m=1,2,3,4$ have types
$A_2,A_3,D_4,\widetilde D_4$; for $m\geq5$ the form is indefinite.

\begin{theorem}[Star family; classical polynomial formula]\label{thm:star}
The endpoints have lengths $m$ and $2m+1$, and
\[
 P_{x,w}(q)=\sum_{j=0}^{\lfloor m/2\rfloor}
 \left(\binom mj-\binom m{j-1}\right)q^j,
 \qquad \binom m{-1}=0.
\]
For $m=2k$ its leading $\mu$-coefficient is
\[
 \mu(x,w)=\binom{2k}{k}-\binom{2k}{k-1}
 =\frac{1}{k+1}\binom{2k}{k}.
\]
For odd $m$ the length difference is even and $\mu(x,w)=0$.
\end{theorem}
The coefficient sequence is unbounded as $k$ grows. This assertion concerns
the changing star systems; only $m=4$ is the affine-$D_4$ case.

\section{The complete interval, not just its rank counts}
\begin{proof}[Proof of Theorem~\ref{thm:star}]
Every $a_Lca_R$ is reduced of length $|L|+|R|+1$. To see this, use the
faithful geometric representation with simple roots $\alpha_c,\alpha_i$.
The initial commuting letters and the subsequent $c$ multiply at ascents.
After $c$, each as-yet unused right letter $a_j$ sees the positive root
\[
 \alpha_c+\sum_{i\in L}\alpha_i+(-1)^{[j\in L]}\alpha_j.
\]
Right multiplication by the other leaves does not change that column.
Thus each right letter is also an ascent.

More generally the image of $\alpha_j$ under $a_Lca_R$ is
\[
 (-1)^{[j\in R]}
 \left(\alpha_c+\sum_{i\in L}\alpha_i+(-1)^{[j\in L]}\alpha_j\right).
\]
Its center coefficient determines whether $j\in R$; its diagonal leaf
coefficient is zero precisely when $j\in L$. Consequently all pairs
$(L,R)$ give distinct elements.

The fixed reduced word for $w$ and the subword criterion now give every
element below it: either a commuting leaf product or an $a_Lca_R$.
The elements at or above $x$ are $x$ itself and exactly those $a_Lca_R$
with $L\cup R=[m]$. For such pairs the order relation is
\[
 a_Lca_R\leq a_{L'}ca_{R'}\iff L\subseteq L',\quad R\subseteq R'.
\]
Indeed, a reduced subword producing an element with $c$-support must
retain $c$, and the distinct-pair formula prevents another representation
from changing either subset. A subword omitting $c$ gives $x$ exactly
when all leaves are selected once.

Each coordinate is therefore left-only, right-only, or both. Order permits
either of the first two states to move to ``both''. With $x$ as the empty
face, this is the entire face lattice of an $m$-cube. There are $3^m+1$
elements. Relative rank $r+1$ contains
$\binom mr2^{m-r}$ elements, for $0\leq r\leq m$.

For a nonempty face $y$, every interval $[y,z]$ is Boolean, so
$R_{y,z}(q)=(q-1)^{\ell(z)-\ell(y)}$ and $P_{y,z}(q)=1$.
For the bottom $x$ and a face $z=a_Lca_R$, we also have
\[
 R_{x,z}(q)=(q-1)^{|L\cap R|+1}.
\]
This follows directly from the ordinary $R$ recurrence: cancel common
right descents $a_R$ from both endpoints, and then cancel common left
descents $a_{[m]\setminus R}$ from both endpoints. The resulting pair is
$(e,a_{L\cap R}c)$, a word with no repeated generator.

Put $F=P_{x,w}$. Ordinary KL reciprocity, using all the nonempty faces,
therefore becomes
\begin{align*}
 q^{m+1}F(q^{-1})-F(q)
 &=\sum_{r=0}^m\binom mr2^{m-r}(q-1)^{r+1}\\
 &=(q-1)(q+1)^m.
\end{align*}
The ordinary degree bound gives $\deg F\leq\lfloor m/2\rfloor$.
In that degree range the coefficient on the left is $-[q^j]F$.
The coefficient on the right is $\binom m{j-1}-\binom mj$, proving
the formula. The asserted $\mu$-values follow from the interval rank $m+1$.
\end{proof}
The Boolean-interval formulas and related lattice formulas are classical
\cite{Brenti}. The top $xcx$ is a boolean reflection: because the leaves
commute it is the palindromic word $a_1\cdots a_mca_m\cdots a_1$.
Both endpoints therefore belong to the older tree-graph boolean-element
framework of Mongelli. His Corollary 2 explicitly displays the star pattern
and its Catalan $\mu$-coefficient; his displayed $f_4$ polynomial is
$1+3q+2q^2$ \cite[preprint pp.6 and 10]{Mongelli}.
The whole proposed family is therefore prior art.

\section{An independently checked affine example}
Use affine-$D_4$ labels with central node $2$ and leaves $0,1,3,4$.
The concrete pair is
\[
 x=s_0s_1s_3s_4,\qquad w=xs_2x,
 \qquad P_{x,w}=1+3q+2q^2,\qquad\mu(x,w)=2.
\]
An independent verifier uses affine signed permutation actions on
$\mathbb R^4$. Generator $0$ is the signed swap of the first two
coordinates, generators $1,2,3$ are adjacent swaps, and generator $4$
is the signed swap of the last two coordinates followed by translation
by $e_3+e_4$. These give the specified Coxeter diagram.

The verifier uses direct subword sets for Bruhat comparison and computes
the polynomials through $R$-reciprocity rather than the discovery program's
KL descent recurrence. It checks all $271$ nontrivial reciprocity identities
in the $272$-element principal lower ideal. The closed interval has $82$
elements with rank counts
\[
 (1,16,32,24,8,1).
\]
The faithful affine actions of the two endpoints, represented as signed
coordinate images followed by translation, are
\[
 x=((-1,-2,-3,-4),(0,0,1,1)),\qquad
 w=((1,3,2,4),(0,-1,1,0)).
\]
All code is saved. Run
\begin{flushleft}\small
\path{python3 research/verify_affine_d4_r.py}\\
Certificate: \path{results/affine-d4-independent.json}\\
Discovery engine: \path{research/broad_cells/targeted.py}
\end{flushleft}

\section{What this establishes, and what it does not}
The affine example disproves the universal statement that a fully
commutative lower endpoint forces $\mu\in\{0,1\}$ in every Coxeter group.
Green's introduction describes a possible universal scope rather than a
numbered universal conjecture \cite[printed pp.2--3]{Green}; the distinction
matters when describing this boundary. We do not claim to be the first to
observe its failure: the explicit polynomial family is in Mongelli's 2011
preprint. The priority report records the exact matching star diagram.

This example does not refute Gern's finite type-$D$ theorem
\cite[Theorem 4.5.11]{Gern} or the companion finite-$E_6/E_7$ theorem.
It also does not, by itself, give a counterexample in every higher affine
type $D$: the degree-four star is not a standard parabolic diagram in
those systems. The family neither requires new combinatorial invariance
nor supplies the sought positive unbounded finite-$D$ model family.

\begin{thebibliography}{9}
\bibitem{Brenti} F. Brenti, A combinatorial formula for Kazhdan--Lusztig
polynomials, Invent. Math. 118 (1994), 371--394, Corollaries 6.8--6.9.
\href{https://fpsac-archive.github.io/FPSAC94/ARTICLES/06Brenti.pdf}{Primary extended abstract}.
\bibitem{Mongelli} P. Mongelli, Kazhdan--Lusztig polynomials of boolean
elements, J. Algebraic Combin. 39 (2014), 497--525.
\href{https://arxiv.org/abs/1111.2945}{arXiv:1111.2945}.
\bibitem{Green} R. M. Green, Leading coefficients of Kazhdan--Lusztig
polynomials and fully commutative elements, J. Algebraic Combin. 30 (2009),
165--171. \href{https://arxiv.org/abs/0801.1650}{arXiv:0801.1650}.
\bibitem{Gern} T. C. Gern, Leading coefficients of Kazhdan--Lusztig
polynomials in type $D$, Ph.D. thesis (2013).
\href{https://arxiv.org/abs/1304.6074}{arXiv:1304.6074}.
\end{thebibliography}
\end{document}
